\documentclass[JEP,XML,SOM,Unicode,NoEqCountersInSection,published]{cedram}
\datereceived{2024-04-08}
\dateaccepted{2025-04-17}
\dateepreuves{2025-05-01}

\usepackage[scr=rsfs,cal=euler]{mathalfa}
\multlinegap0pt
\let\epsilon\varepsilon

\hyphenation{authors rela-tion restric-tion}

\let\ts\textstyle
\newcommand\mto{\mathchoice{\longmapsto}{\mapsto}{\mapsto}{\mapsto}}

\newcommand{\Psfrac}[2]{(\sfrac{#1}{#2})}
\newcommand{\psfrac}[2]{\sfrac{(#1)}{#2}}
\newcommand{\spfrac}[2]{\sfrac{#1}{(#2)}}
\newcommand{\pspfrac}[2]{\sfrac{(#1)}{(#2)}}

\newcommand{\app}{\mathrm{app}}
\newcommand{\loc}{\mathrm{loc}}
\newcommand{\st}{\mathrm{st}}

\def\Xint#1{\mathchoice
{\XXint\displaystyle\textstyle{#1}}%
{\XXint\textstyle\scriptstyle{#1}}%
{\XXint\scriptstyle\scriptscriptstyle{#1}}%
{\XXint\scriptscriptstyle\scriptscriptstyle{#1}}%
\!\int}
\def\XXint#1#2#3{{\setbox0=\hbox{$#1{#2#3}{\int}$ }
\vcenter{\hbox{$#2#3$ }}\kern-.6\wd0}}
\def\ddashint{\Xint=}
\def\dashint{\Xint-}

\renewcommand{\div}{\operatorname{div}}

\theoremstyle{plain}
\newtheorem{thm}{Theorem}[section]
\newtheorem{prop}[thm]{Proposition}
\newtheorem{coro}[thm]{Corollary}
\newtheorem{lemme}[thm]{Lemma}

\theoremstyle{definition}
\newtheorem{defi}[thm]{Definition}
\newtheorem{rem}{Remark}[section]

\let\pa\partial
\let\na\nabla
\let\eps\varepsilon

\newcommand{\eff}{\mathrm{eff}}
\newcommand{\sym}{\operatorname{sym}}
\newcommand{\ssym}{\operatorname{ssym}}
\newcommand{\Lip}{\mathrm{Lip}}
\newcommand{\supp}{\operatorname{supp}}
\newcommand{\mO}{\mathcal{O}}
\newcommand{\mM}{\mathcal{M}}
\newcommand{\mS}{\mathcal{S}}
\newcommand{\mV}{\mathcal{V}}
\newcommand{\R}{\mathbb{R}}
\newcommand{\T}{\mathbb{T}}
\newcommand{\N}{\mathbb{N}}
\newcommand{\Z}{\mathbb{Z}}

\datepublished{2025-05-13}
\begin{document}
\frontmatter
\title{Analysis of a sedimenting suspension near a\nobreakspace vertical wall}

\author[\initial{D.} \lastname{Gérard-Varet}]{\firstname{David} \lastname{Gérard-Varet}}
\address{Université Paris Cité and Sorbonne Université, CNRS, IMJ-PRG,\\ F-75013 Paris, France}
\email{david.gerard-varet@imj-prg.fr}
\urladdr{https://dgerardv.github.io/}

\author[\initial{A.} \lastname{Mecherbet}]{\firstname{Amina} \lastname{Mecherbet}}
\address{Université Paris Cité and Sorbonne Université, CNRS, IMJ-PRG,\\ F-75013 Paris, France}
\email{mecherbet@imj-prg.fr}
\urladdr{https://webusers.imj-prg.fr/~amina.mecherbet/}

\thanks{The work of D.\,G-V was supported by project Singflows grant ANR-18-CE-40-0027 and project Bourgeons grant ANR-23-CE40-0014-01 of the French National Research Agency (ANR). Both authors were partially supported by the Agence Nationale de la Recherche, Project SUSPENSIONS, grant ANR-24-CE92-0028-01}

\begin{abstract}
We consider a sedimenting suspension in a Stokes flow, in the presence of a vertical wall. We study the effect of a particle-depleted fluid layer near the wall on the bulk dynamics of the suspension. We show that this effect can be captured by an appropriate wall law of Navier type. Under appropriate lower bound for the minimal distance and the size of the depletion layer, we~provide in this way a rigorous justification of the apparent slip observed in many experiments. We also discuss the phenomenon of intrinsic convection predicted in some physics articles.
\end{abstract}

\subjclass{35Q35, 76T20, 76D07, 35Q70, 76D10}

\keywords{Sedimentation of particles, Stokes flow, method of reflections, boundary layer analysis}

\altkeywords{Sédimentation de particules, fluide de Stokes, méthode de réflexions, analyse de type couche limite}

\alttitle{Analyse d'une suspension de particules en sédimentation à proximité d'une paroi verticale}

\begin{altabstract}
Nous considérons des particules sédimentant dans un fluide de Stokes en la présence d'une paroi verticale. Nous étudions l'effet du bord ainsi que de la couche limite appauvrie en particules sur la dynamique de la suspension. Nous montrons que cet effet peut être modélisé par des conditions au bord de type Navier. Sous certaines hypothèses concernant la distance minimale entre particules ainsi que la taille de la couche limite, nous proposons une justification rigoureuse du phénomène de glissement apparent observé dans certaines expériences. Nous discutons également du phénomène de convection intrinsèque étudié dans certains articles de physique.
\end{altabstract}

\maketitle

\tableofcontents
\mainmatter

\section{Introduction}
The mathematical analysis of suspensions of small particles sedimenting in a Stokes flow has been the subject of extensive research over the last years. In~the context of inertialess particles, a classical model is the following. Given a finite or infinite collection of identical non-overlapping spherical particles $(B_i = B(X_i,R))_{i \in I}$, set in a bounded or unbounded domain $\Omega$, one considers the system
\begin{equation} \label{stokes_equation}
\begin{aligned}
-\mu \Delta u + \na p & = f \quad \text{ in } \: \Omega \setminus(\ts\bigcup_i B_i), \\
\div u & = 0 \quad \text{ in } \: \Omega \setminus(\ts\bigcup_i B_i), \\
u\vert_{\pa \Omega} & = 0, \\
D(u) & = 0 \quad \text{ in } B_i \quad \forall i, \\
\int_{\pa B_i} \Sigma_\mu(u,p) n d \sigma & = -m g e, \quad \int_{\pa B_i} [\Sigma_\mu(u,p) n] \times (x-X_i) d \sigma(x) = 0, \quad \forall i.
\end{aligned}
\end{equation}
The unknowns of this system are the velocity field and pressure:
\[
u(x) = (u_1,u_2,u_3)(x_1,x_2,x_3), \quad p(x) = p(x_1,x_2,x_3).
\]
The first three lines describe a Stokes flow with viscosity $\mu$ outside the particles, with the usual no-slip condition at $\pa \Omega$. The fourth one expresses rigidity of the velocity field of each particle $B_i$, and is equivalent to
\[
u = v_i + \omega_i \times (x-X_i) \quad \text{in } B_i
\]
for some unknown translational and angular velocities $v_i, \omega_i \in \R^3$. The two relations in the last line, correspond to balances of forces and torques, and involve the Newtonian stress tensor $\Sigma_\mu(u,p) = 2\mu D(u) - p I_3$ with $D(u)=\frac{1}{2}(\nabla u + \nabla^\top u)$ the symmetric gradient of $u$. The constants $v_i, \omega_i$ can be seen as Lagrange multipliers associated with these constraints. The term $-mge$, $e = (0,0,1)^t$ encodes gravity and buoyancy effects: denoting $\rho_p$ and $\rho_f$ the particle and fluid masses per unit volume, one has $m = \frac{4}{3} \pi R^3 (\rho_p - \rho_f)$, while $g$ is the usual gravitational constant. This is the term responsible for sedimentation.

Roughly, all mathematical studies of this model divide into two categories:
\begin{itemize}
\item Those that focus on steady properties of the suspension. Time evolution of $X_i = X_i(t)$ is neglected (or frozen), the particle distribution is considered as given, and mean properties of the suspension, such as the average settling velocity \cite{HillairetHoefer23, Hasimoto59, Batchelor72Sedimentation, DuerinckxGloria22} or the effective viscosity \cite{HainesMazzucato12,NiethammerSchubert19, Gerard-Varet19, Gerard-VaretHillairet19, Gerard-VaretHoefer21, Gerard-VaretMecherbet20,DuerinckxGloria19}, are deduced. The derivation of the mean settling velocity, and its expansion in terms of the solid volume fraction are notoriously difficult, and rely on subtle cancellation mechanisms which require strong assumptions on the distribution of the particles (either periodic-like or random stationary with strong mixing properties).
\item Those that analyze the dynamics of the suspension (for a large but finite $N \gg 1$ collection of balls $B_1, \dots, B_N$): the time dependence of $X_i = X_i(t)$ is taken into account through coupling of \eqref{stokes_equation} with the ODE
\begin{equation}
\dot{X}_i = v_i.
\end{equation}
In the dilute case, under appropriate lower bound on the minimal distance between the particles, a mean field analysis allows to derive a transport-Stokes equation on
\[
\rho = \lim_{N\rightarrow \infty} \frac{1}{N} \sum_{i=1}^N \delta_{X_i}.
\]
We refer to \cite{Hofer18MeanField,Mecherbet19,Hofer&Schubert,HoferSchubert23} for all details.
\end{itemize}

Except for the work \cite{HillairetHoefer23} that contains a brief discussion of boundary effects, none of the mathematical references above investigate the role of the boundaries on the dynamics of the suspension.

On the contrary, the impact of the wall of a container on the nature of the sedimentation process has been widely discussed in the physics community. One of the most striking examples of such impact is the so-called Boycott effect \cite{Boy20}: the settling speed of a suspension is much faster in an inclined channel than in a vertical one, as the slurry of particles accumulating on the upward-facing wall slides faster to the bottom, creating a counterflow in the clear fluid layer above.
Even in a vertical channel, the role of boundaries has been emphasized for decades. One main reason for this boundary effect is the formation of a so-called depletion layer near the wall, where the concentration of the suspension is much less than in the bulk of the suspension, and where large velocity gradients occur. This depletion may be due to excluded volume effect: for instance, for spherical particles, one has the constraint $d(X_i, \pa \Omega) > R$, which results in a poorer concentration in a layer of thickness $\approx R$. But additional mechanisms, such as shear induced migration, may enhance the depletion mechanism and much enlarge the layer \cite{SegSil,SegSil2,Ho_Leal_1974,Ausserre}, notably for suspensions of polymers. In~relation to this depletion layer, two phenomena are much discussed in the literature:
\begin{itemize}
\item The first one, observed experimentally, is {\em apparent slip}: for suspensions flowing through a channel, one observes non-zero downward velocities just outside the depletion layer, that can be interpreted as additional slip : see \cite{Abbasi,Koponen,Ghosh} and references therein.
\item The second one, not observed experimentally but predicted in several studies \cite{GeiMaz,BFAH96,BFBA98}, is called {\em intrinsic convection}: homogeneous suspensions are expected to experience an additional downward flow in the bulk, while experiencing upward counterflows just outside the depletion layers (resulting in zero additional net flux).
\end{itemize}
At first sight, these two phenomena may seem contradictory (downward {\it vs} upward flow just outside the depletion layers). One ambition of this paper is to clarify these different behaviors. More globally, we~wish here to analyze rigorously the effect of a depletion layer on a suspension, starting from the simplest model \eqref{stokes_equation}. For $\eps > 0$ we consider the half-space $\Omega = \Omega^\eps = (-\eps,+\infty) \times \R^2$, where the depletion layer $D_\eps = (-\eps,0) \times \R^2$, of width $\eps$, is free of particles centers. Namely,
\begin{equation} \label{eq_assumption}
X_1, \dots, X_N \in K \Subset \overline{\Omega^0}, \quad \Omega^0 = \R^*_+ \times \R^2.
\end{equation}
Taking
\[
L = |K|^{1/3}, \quad U = \frac{mgN}{\mu L}, \quad P = \frac{\mu U}{L},
\]
as reference length, velocity and pressure, we~can put system \eqref{stokes_equation} in dimensionless form:
\begin{align*}
& X' = X/L, \quad \eps' = \eps/L, \quad R'= R/L, \\
& u' = u/U, \quad p' = p/P, \quad f' = \frac{L^2}{\mu U} f, \quad X'_i = X_i/L, \quad B'_i = B_i/L = B(X'_i, R').
\end{align*}
After dropping the primes, we~end up with
\begin{equation} \label{main_equation}
\begin{aligned}
-\Delta u + \na p & = f \quad \text{ in } \: \Omega^\eps \setminus(\bigcup B_i), \\
\div u & = 0 \quad \text{ in } \: \Omega ^\eps\setminus(\bigcup B_i), \\
u\vert_{\pa \Omega^\eps} & = 0, \\
D(u) & = 0 \quad \text{ in } B_i, \quad 1 \le i \le N, \\
\int_{\pa B_i} \Sigma(u,p) n d \sigma & = -\frac{1}{N}e, \quad \int_{\pa B_i} \Sigma(u,p) n \times (x-X_i) d \sigma(x)= 0, \quad 1 \le i \le N,
\end{aligned}
\end{equation}
where $\Sigma(u,p) = 2D(u) - p I$. Note that we could have considered some variation of~\eqref{main_equation} with additional forcing $\int_{B_i} f$ and torque $\int_{B_i} f \times (x-X_i)$ included in the last two relations.
We want to understand the asymptotic behaviour of the solution $u = u^{N,\eps}$ for large $N$ and small $\eps$. We make three main assumptions :
\begin{itemize}
\item We denote by $\rho^N$ the empirical measure defined by
$
\rho^N = \frac{1}{N} \sum_i \delta_{X_i}
$
and consider a bounded density~$\rho$ compactly supported in $\overline{\Omega^0}$.
\item We assume that there exists $\theta>1$ such that
\begin{equation} \label{H1}
\epsilon> \theta R,
\end{equation}
which is a slight reinforcement of the non-overlapping constraint between the spheres and the wall.
\item Denoting $d_{\min} = \min_{i \neq j} |X_i - X_j|$, we~assume that
\begin{equation}\label{hyp_dmin}
d_{\min}\geq \max\Bigl (\frac{C}{N^{1/3}},2\theta R \Bigr).
\end{equation}
\end{itemize}
\begin{rem}
The three assumptions above, notably the lower bound \eqref{H1} on the size of the depletion layer, are enough to prove our Theorems \ref{thm1} and \ref{thm2} and Corollary~\ref{coro1}. But, as will be discussed in Section \ref{subsec_apparent}, in order to relate these mathematical results to the apparent slip phenomenon, one needs a more stringent condition than \eqref{H1}, essentially
\[
\epsilon \ge \mathcal{C} (R+ \rho+ \|\rho^N-\rho\|),
\]
where $\mathcal{C}$ is a large enough constant independent of $N$ and $R$, and where $\| \, \|$ is a weak norm, see \eqref{stronger_assumption_size}. This stronger hypothesis will be discussed in detail in Section \ref{sec_final_comments} as well as the phenomenon of intrinsic convection.
\end{rem}
Under these assumptions, we~will perform a two-step analysis of \eqref{main_equation}.
\subsubsection*{Step 1 : approximation by a continuous model.}
We will first show that the solution $u^{N,\eps}$ of \eqref{main_equation} can be approximated by the continuous solution $u^\eps$ of
\begin{equation} \label{limit_equation}
\begin{aligned}
-\Delta u^\eps + \na p^\eps & = f - \rho e \quad \text{ in } \: \Omega^\eps, \\
\div u^\eps & = 0 \quad \text{ in } \: \Omega^\eps, \\
u^\eps \vert_{\pa \Omega^\eps} & = 0.
\end{aligned}
\end{equation}
We assume in all what follows that $f \in L^1(\Omega^\eps) \cap L^\infty(\Omega^\eps)$. In~particular, $f \in L^{6/5}(\Omega^\eps)$ so that \eqref{main_equation}, \resp \eqref{limit_equation}, has a unique solution $u^{N,\eps}$, \resp $u^\eps$ in
\[
V = \{ u \in H^1_{\loc}(\overline{\Omega^\eps}), \na u \in L^2(\Omega^\eps)\}.
\]
We will prove in Section \ref{sec_approx}:
\begin{thm} \label{thm1}
Let $1<q<3/2$, $Q \Subset \overline{\Omega^\eps}$. There exists $C > 0$ depending on $q$, $Q$ and on the constant $\theta$ in \eqref{H1}, such that ,
\begin{equation} \label{first_estimate}
\|u^{N,\eps} - u^\eps\|_{L^q(Q)} \le C \big(\phi + \|\rho^N-\rho\|_{(W^{2,q'}(\Omega^0))^*}+R \big),
\end{equation}
where $\phi=R^3N$ the particles volume fraction.
\end{thm}
\begin{itemize}
\item The norm $\|\rho^N-\rho\|_{(W^{2,q'}(\Omega^0))^*}$ can be replaced with $W_1(\rho^N,\rho)$, the first Wasserstein distance between $\rho^N$ and $\rho$, and consequently by $W_p(\rho^N,\rho)$ for any $p\in[1,+\infty]$ if one assumes $\rho \in \mathcal{P}(\Omega^0)$. See Remark \ref{rem:Wasserstein_versus_sob_neg}.
\item The control on the minimal distance \eqref{hyp_dmin} can be relaxed to $d_{\min}>2 \theta R $ together with any assumption ensuring a uniform bound on
\[
\frac{1}{N} \sum_{j \neq i} \frac{1 }{|X_i-X_j|^2} \leq C,
\]
see Remark \ref{discussion_d_min} for more details.
\end{itemize}

In cases without boundaries, $\Omega = \R^3$ or $\Omega = \T^3$, this kind of statement is already known, see for instance \cite{Hofer&Schubert}. The proof is based there on an approximation of $u^{N,\eps}$ made of so-called point Stokeslets and point Stresslets, centered at the $X_i$'s. They involve the Oseen tensor, that is the fundamental solution of the Stokes operator in~$\R^3$, and its derivatives. They are not fitted to our half-space case, as the inhomogeneous Dirichlet data that they create at the boundary is too large. We must therefore adapt this strategy, constructing the analogue of point Stokeslets for a half-space with homogeneous Dirichlet condition. We believe that this part of the analysis is of independent interest, and contributes to the interesting problem of extending the method of reflections to domains with boundaries \cite{Hoefer19}.

\subsubsection*{Step 2 : wall law.}
Once the continuous model \eqref{limit_equation} has been derived, the goal is to understand the effect of the depletion layer $D^\eps$ on the solution $u^\eps$. In~other words, one would like to improve the crude approximation
\begin{equation} \label{crude_approx}
\begin{aligned}
-\Delta u^0 + \na p^0 &= f - \rho e \quad \text{ in } \: \Omega, \\
\div u^0 & = 0 \quad \text{ in } \: \Omega^0, \\
u^0 \vert_{\pa \Omega^0} & = 0,
\end{aligned}
\end{equation}
where the impact of the layer is totally neglected. In~this regard, we~will prove in Section \ref{sec_wall_law} the following theorem:
\begin{thm} \label{thm2}
Assume that $f$ and $\rho$ are smooth, compactly supported in $\overline{\Omega^0}$. Then, for all $m \in \N$, one has
\[
u^\eps = u^0 + \eps u^1 + \dots + \eps^m u^m + O(\eps^{m+1}) \quad \text{ in } \: \dot{H}^1(\Omega^0),
\]
where $u^0$ solves \eqref{crude_approx}, and for each $i \ge 1$, $u^i$ solves an homogeneous Stokes equation in~$\Omega^0$ with inhomogeneous Dirichlet data at $\pa \Omega^0$ coming from lower order profiles.
\end{thm}
This theorem provides an asymptotic expansion of $u^\eps$ at arbitrary order, where each successive correction refines the description of the effect of the depletion layer. One could relax the assumptions on the smoothness of $\rho$ and $f$: for instance, it is enough that $\na u^0,p^0$ belong to $H^{m+2}(\Omega^0)$ to have an expansion with $m$ terms.

The analysis leading to Theorem \ref{thm2} shows in particular that the inhomogeneous Dirichlet data for $u^1$ is given by
\begin{equation} \label{corrector_equation}
u^1 \vert_{\pa \Omega^0} = (0, \pa_1 u^0_{2},\pa_1 u^0_{3})\vert_{\pa \Omega^0}.
\end{equation}
This will imply, see Section \ref{sec_wall_law}:
\begin{coro} \label{coro1}
Let $\rho$ and $f$ as in the previous theorem, and $u^S$ the solution of the Stokes system with Navier slip boundary condition:
\begin{equation} \label{improved_approximation}
\begin{aligned}
-\Delta u^S + \na p^S & = f - \rho e \quad \text{ in } \: \Omega^0, \\
\div u^S & = 0 \quad \text{ in } \: \Omega^0, \\
u^S\vert_{\pa \Omega^0} & = \eps (0, \pa_1 u^S_2, \pa_1 u^S_3)\vert_{\pa \Omega^0}.
\end{aligned}
\end{equation}
Then, one has the estimate
\[
\|u^\eps-u^S\|_{\dot{H}^1(\Omega^0)} = O(\eps^2).
\]
\end{coro}
This $O(\eps^2)$ estimate shows that imposing a wall law of Navier type at the artificial boundary
$\pa \Omega^0$ improves the $O(\eps)$ error estimate provided by the homogeneous Dirichlet condition. It is the rigorous translation of the notion of apparent slip evoked in several papers, \cite{Abbasi,Koponen, Ghosh}. We will further discuss it in Section \ref{sec_final_comments}. The other phenomenon mentioned above, that is {\em intrinsic convection} is more subtle and corresponds to a degenerate regime of what we consider here. We will also discuss it in Section \ref{sec_final_comments}.

\subsubsection*{Acknowledgements}
The authors would like to thank Matthieu Hillairet and Richard Höfer for stimulating discussions on the topic.

\section{Stokes problem on the half space: fundamental solution and Stokeslets} \label{sec_stokes}
For the analysis carried in this section \ref{sec_stokes} and the next section \ref{sec_approx}, there is no restriction in taking $\eps = 0$, that is working on $\Omega^0 =(0,+\infty)\times \R^2$. The results adapt straightforwardly to $\Omega^\epsilon=(-\varepsilon,+\infty)\times \R^2$ by translation.

We introduce the fundamental solution to the Stokes equation on the half space $(0,+\infty)\times \R^2$. Given $y \in (0,+\infty)\times \R^2$, we~set $x \mto G(x,y)$ the unique solution to
\begin{equation}
\left\{
\begin{array}{rcl}
-\Delta u+\nabla q &=& \delta_{y} \mathbb{I} \quad\text{on } \Omega^0, \\
\div(u)&=&0 \quad\text{on } \Omega^0, \\
u\vert_{x_1=0}&=&0,
\end{array}
\right.
\end{equation}
where $\mathbb{I}$ is the identity matrix in $3d$. We have for all $x \neq y$
\begin{equation}\label{ineq:fundamental_solution}
|G(x,y)| \leq \frac{C}{|x-y|}, \quad |\nabla_y G(x,y)|+ |\nabla_x G(x,y)| \leq \frac{C}{|x-y|^2}.
\end{equation}
The proof of such an estimate is postponed to the appendix.
Let us also introduce the Stokes velocity for a sedimenting sphere in a half space. Given $F\in \R^3$, $r>0$ and $y_0 \in \Omega^0$ such that $(y_0)_1>r$, we~use the notation $x \mto \mathcal{U}^{\st}_{r}[F](x,y_0)\in \dot{H}^1(\Omega^0)$ for the unique solution of the Stokes equation
\begin{gather} \label{eq_stokeslet}
\left\{
\begin{array}{rcl}
-\Delta u+\nabla q &=& 0\quad\text{on } \Omega^0 \setminus \overline{B(y_0,r)},\\
\div(u)&=&0 \quad\text{on }\Omega^0 \setminus \overline{B(y_0,r)},\\
D(u) &=&0 \text { on } B(y_0,r),\\
u\vert_{x_1=0}&=&0,
\end{array}
\right.
\\
\label{eq_stokeslet: force_torque}
\int_{\partial B(y_0,r)} \Sigma(u,p) n d \sigma =F,\quad
\int_{\partial B(y_0,r)} [\Sigma(u,p) n] \times(x-y_0)d \sigma(x)=0.
\end{gather}
We denote by $\mathcal{P}^{\st}_{r}[F](x,y_0) $ its associated pressure. Note that $\mathcal{U}^{\st}_{r}[F](\cdot,y_0)$ is implicitly given by a rigid vector field in $B(y_0,r)$. The associated pressure will be denoted by $\mathcal{P}^{\st}_{r}[F](\cdot,y_0)$. We have the following
\begin{prop}\label{prop:stokeslet}
Let $r>0$, $\theta>1$ and $y_0 \in \Omega_0$ such that $ (y_0)_1> \theta r$.
\begin{enumerate}
\item We have, for all $x \in \Omega^0 \setminus B(y_0,r)$,\label{ref:it1}
\[
\mathcal{U}^{\st}_{r}[F](x,y_0)= \frac{1}{r}\, \mathcal{U}^{\st}_{1}[F](\sfrac{x}{r},\sfrac{y_0}{r}).
\]

\item There exists $\mathcal{H}_r[F](\cdot,y_0)$ such that, for all $ x \not \in B(y_0,\theta r)$, \label{ref:it2}
\begin{gather}\label{eq:decomposition_stokeslet}
\mathcal{U}^{\st}_{r}[F](x,y_0)= G(x,y_0)F + \mathcal{H}_r[F](x,y_0),
\\
\label{eq:decay_rest_stokeslet}
\left|\mathcal{H}_r[F](x,y_0) \right|\leq r \frac{C_\theta|F| }{|x-y_0|^2}, \quad \left|\nabla_x \mathcal{H}_r[F](x,y_0) \right|\leq r \frac{C_\theta |F|}{|x-y_0|^3}.
\end{gather}
\item The energy satisfies \label{ref:it3}
\[
\|\nabla \mathcal{U}^{\st}_r[F](\cdot,y_0)\|_{L^2(\Omega^0)}^2=2\|D \left(\mathcal{U}^{\st}_r[F](\cdot,y_0) \right) \|_{L^2(\Omega^0)}^2
\]
and
\begin{equation}\label{eq:energy}
\|\nabla \mathcal{U}^{\st}_r[F](\cdot,y_0)\|_{L^2(\Omega^0)}\leq \frac{C|F|}{\sqrt{r}},
\end{equation}
with a constant $C>0$ independent of $y_0$, $F$ and $r$.
\end{enumerate}
\end{prop}
In order to prove such a result we first recall the following extension result which can be obtained by standard scaling arguments.
\begin{lemme}\label{lemme:extension}
Let $p>1$, $\theta>1$, $a \in \R^3$, $\lambda>0$ and $v \in W^{1,p}(B(a,\lambda))$ a divergence-free velocity field with vanishing mean on $B(a,\lambda)$. There exists an extension $u_v \in W^{1,p}_0(B(a,\theta \lambda))$
of $v$ which is divergence-free and such that
\[
\|\na u_v\|_{L^p(B(a,\theta \lambda) \setminus B(a,\lambda))} \leq C_\theta \|\nabla v\|_{L^p(B(a,\lambda))},
\]
with a constant depending only on $\theta$ and $p$ and not on $\lambda$ and $a$.
\end{lemme}

\skpt
\begin{proof}[Proof of Proposition \ref{prop:stokeslet}]
\eqref{ref:it1} The proof relies on a standard scaling argument.

\eqref{ref:it3}
The first identity is valid for any divergence-free vector field $v$ vanishing at $x_1 =\nobreak 0$: it follows from the identity $-\Delta v = -2 \div(Dv)$, multiplying by $v$ and integrating by parts. As regards the energy bound, we~use the following:
\begin{align*}
\|\nabla \mathcal{U}^{\st}_1[F]&(\cdot,y_0)\|_{L^2((0,+\infty)\times \R^2)}^2=2\|D\left(\mathcal{U}^{\st}_1[F](\cdot,y_0)\right) \|_{L^2((0,+\infty)\times \R^2)}^2\\
&=2\|D \left(\mathcal{U}^{\st}_1[F](\cdot,y_0) \right) \|_{L^2((0,+\infty)\times \R^2 \setminus B(y_0,1))}^2\\
&= \int_{(0,+\infty)\times \R^2 \setminus B(y_0,1)}\!\! \left(2D\left(\mathcal{U}^{\st}_1[F](\cdot,y_0)\right)- \mathcal{P}^{\st}_1[F](\cdot,y_0) I_3\right) : D\left(\mathcal{U}^{\st}_1[F](\cdot,y_0)\right)\\
&= \int_{(0,+\infty)\times \R^2 \setminus B(y_0,1)}\!\! \left(- \Delta \mathcal{U}^{\st}_1[F](\cdot,y_0)+ \nabla \mathcal{P}^{\st}_1[F](\cdot,y_0)\right) \cdot \mathcal{U}^{\st}_1[F](\cdot,y_0)\\
&\hspace*{2cm}+
\int_{\partial B(y_0,1)} \Sigma\left (\mathcal{U}^{\st}_1[F](\cdot,y_0),\mathcal{P}^{\st}_1[F](\cdot,y_0)\right)n \cdot \mathcal{U}^{\st}_1[F](\cdot,y_0) d \sigma\\
&=F \cdot \dashint_{B(y_0,1)} \mathcal{U}^{\st}_1[F](\cdot,y_0) \\
&\leq C |F| \| \nabla \mathcal{U}^{\st}_1[F](\cdot,y_0) \|_{L^2}.
\end{align*}
Note that, for the last equality we used the fact that $\mathcal{U}^{\st}_1[F](\cdot,y_0)$ is a rigid motion on $B(y_0,1)$. For the last inequality we have used the Sobolev embedding
\[
\{ u \in \dot{H}^1(\Omega^0), \; u\vert_{\pa \Omega^0}= 0 \} \subset L^6(\Omega^0).
\]
For general $r$, we~then use the scaling argument of \eqref{ref:it1}.

\eqref{ref:it2} We need to show that\vspace*{-3pt}
\[
\mathcal{H}_r[F](\cdot,y_0) :=\mathcal{U}^{\st}_{r}[F](\cdot,y_0) - G(\cdot,y_0) F
\]
satisfies the inequalities in \eqref{eq:decay_rest_stokeslet}. Note that $\frac{1}{r}G(\sfrac{x}{r},\sfrac{y_0}{r})\!=\!G(x,y_0)$. Together with~\eqref{ref:it1}, it~shows that $\mathcal{H}_r[F]$ inherits the scaling property:
$\mathcal{H}_{r}[F](x,y_0)\!=\! \frac{1}{r} \mathcal{H}_{1}[F](\sfrac{x}{r},\sfrac{y_0}{r})$.
This allows to restrict to $r=1$. We have for all $x \in \Omega^0 \setminus B(y_0,\theta)$\vspace*{-5pt}
\begin{multline*}
\mathcal{U}^{\st}_{1}[F](x,y_0) - \biggl(\dashint_{ B(y_0,1)} G(x,z) d z \biggr) F = \\
\int_{\partial B(y_0,1)} \biggl(G(x,y)-\dashint_{ B(y_0,1)} G(x,z) d z \biggr)\Sigma\left(\mathcal{U}^{\st}_{1}[F](y,y_0),\mathcal{P}^{\st}_1[F](\cdot,y_0)\right) n d \sigma (y).
\end{multline*}
Hence, using Lemma \ref{lemme:extension}, we~denote by $v\in H^1_0(B(y_0,\theta))$ a divergence-free extension of $y \mto G(x,y)-\dashint_{ B(y_0,1)} G(x,z) d z$ such that
\[
\| \nabla v \|_{L^2(B(y_0,\theta))} \leq C_\theta \|y\mapsto \nabla_y G(x,y) \|_{L^2(B(y_0,1))},
\]
with a constant depending only on $\theta$, and one gets\vspace*{-3pt}
\begin{multline*}
\biggl |\mathcal{U}^{\st}_{1}[F](x,y_0)- \biggl(\dashint_{ B(y_0,1)} G(x,z) d z \biggr) F \biggr|= \biggl |\int_{B(y_0,\theta)\setminus B(y_0,1)} \nabla v : \nabla \mathcal{U}^{\st}_{1}[F](\cdot,y_0) \biggr|\\
\leq \| \nabla \mathcal{U}^{\st}_{1}[F](\cdot,y_0)\|_{L^2} \| \nabla v \|_{L^2(B(y_0,\theta))}.
\end{multline*}
We then use the energy estimate for $ \mathcal{U}^{\st}_{1}[F]$ of item \eqref{ref:it3}
together with inequality \eqref{ineq:fundamental_solution} which yields\vspace*{-3pt}
\begin{equation}\label{ineq1}
\| \nabla v \|_{L^2(B(y_0,\theta))} \le C_\theta \|y\mto \nabla_y G(x,y)\|_{L^\infty(B(y_0,1))} \leq \frac{C'_\theta}{|x-y_0|^2},
\end{equation}
where we used that $|x-y|>(1-1/\theta)|x-y_0|$ provided that $|x-y_0|>\theta$. Hence, we~write\vspace*{-3pt}
\[
\mathcal{H}_1[F](\cdot,y_0) = \mathcal{U}^{\st}_{1}[F](\cdot,y_0) - \biggl(\dashint_{ B(y_0,1)}\hspace*{-2mm} G(\cdot,z) dz \biggr) F
+ \biggl(\dashint_{ B(y_0,1)}\hspace*{-1mm} [G(\cdot,z) - G(\cdot,y_0)] dz \biggr) F
\]
and the extra term $\bigl(\dashint_{ B(y_0,1)} [G(\cdot,z) - G(\cdot,y_0)] dz \bigr) F$ can be estimated analogously using~\eqref{ineq1}. This concludes the proof of the first inequality in \eqref{eq:decay_rest_stokeslet}. The second one is similar, applying $\na_x$ to all quantities above, and using the second bound in \eqref{ineq:fundamental_solution} instead of the first one.
\end{proof}

\section{Approximation of the velocity field} \label{sec_approx}

The main goal of this section is to prove Theorem \ref{thm1}. There is no loss of generality in proving it in the case $\eps = 0$. We use the shortcuts $u^N = u^{N,0}$, $F=-\Psfrac{1}{N}e $. Let us denote\vspace*{-3pt}
\[
u_{\app}= \sum_i \mathcal{U}^{\st}_R\left[F\right](x,X_i)+ u_{f,N},
\]
where we recall that $\mathcal{U}^{\st}_R[F](\cdot,X_i)$ is the Stokeslet attached to $B_i=B(X_i,R)$, see~\eqref{eq_stokeslet}, and where $u_{f,N}$ is the solution of the Stokes equation\vspace*{-3pt}
\begin{equation}\label{eq:u_f}
\begin{cases}
\begin{aligned}
-\Delta u_{f,N}+\nabla q_{f,N} &= f (1-1_{\oldbigcup B_i}) \quad\text{on }(0,+\infty)\times \R^2, \\
\div(u_{f,N})&=0 \quad\text{on }(0,+\infty)\times \R^2, \\
u_{f,N}\vert_{x_1=0}&=0.
\end{aligned}
\end{cases}
\end{equation}
We first aim to estimate the error between $u^N$ and
$u_{\app}$.

\Subsection{From $u^N$ to $u_{\app}$}
\begin{prop}
Let $K \Subset \overline{\Omega^0}$. We have, for any $q<3/2$,\vspace*{-3pt}
\[
\|u^N-u_{\app}\|_{L^q(K)}\leq C_{K} \phi.
\]
\end{prop}
\begin{proof}
We set $v=u^N-u_{\app}$, it satisfies\vspace*{-3pt}
\begin{gather*}
\begin{cases}
\begin{aligned}
-\Delta v+\nabla q &= 0 \quad\text{on }\Omega^0 \setminus \overline{\bigcup B_i},\\
\div(v)&=0 \quad\text{on }\Omega^0 \setminus \overline{\bigcup B_i},\\
D(v) &=- D(u_{\app}) \quad\text {on } B_i,\\
v\vert_{x_1=0}&=0,
\end{aligned}
\end{cases}
\\
\int_{B_i} \Sigma(v,q) n =
\int_{B_i} [\Sigma(v,q) n] \times(x-X_i)=0.
\end{gather*}

In order to estimate the $L^p_{\loc}$ norm of $v$, we~will use a duality argument, observing that for any $K\Subset \overline{\Omega^0}$ and $\psi \in C^\infty_c(K)$
\[
\int_{K} v \cdot \psi = -2 \int_{\Omega^0} D (v): D (u_\psi),
\]
with $u_\psi \in W^{2,q'}(\Omega^0) \cap W^{1,q'}_0(\Omega^0)$ being the unique solution to the Stokes equation $-\Delta u_\psi+\nabla q_\psi=\psi $, $\div u_\psi=0$ on $\Omega^0$ with vanishing Dirichlet boundary condition on $x_1=0$. We claim that this solution satisfies
\begin{equation}\label{reg_u_psi}
\|u_\psi \|_{W^{2,q'}(\Omega^0)} \leq C_K \|\psi\|_{L^{q'}(K)}.
\end{equation}
Indeed, the bound
\[
\|\na^2 u_\psi\|_{L^{q'}(\Omega^0)} \le C \|\psi\|_{L^{q'}(\Omega^0)} = C \|\psi\|_{L^{q'}(K)}
\]
follows from \cite[Th.\,IV.3.2,\,(IV.3.30)]{Galdi}. Moreover, by taking $\tilde{q} \!\in\! (1,3)$ such that $\spfrac{3 \tilde q}{3-\tilde q} = q'$, we~deduce from \cite[Th.\,IV.3.2,\,(IV.3.31)]{Galdi} that
\[
\|\na u_\psi\|_{L^{q'}(\Omega_0)} \le C \|\psi\|_{L^{\tilde q}(\Omega^0)} = C \|\psi\|_{L^{\tilde q}(K)} \le C_K \|\psi\|_{L^{q'}(K)}.
\]
Finally, the bound
\[
\|u_\psi\|_{L^{q'}(\Omega_0)} \le C \|\psi\|_{L^{q'}(K)}
\]
is shown as follows. Let $R > 0$ such that $K \subset B(0,R-1)$. From Poincaré inequality (remember that $u_\psi$ vanishes at $\pa \Omega_0$):
\[
\|u_\psi\|_{L^{q'}(\Omega_0 \cap B(0,R))} \le C \|\na u_\psi\|_{L^{q'}(\Omega_0 \cap B(0,R))} \le C_K \|\psi\|_{L^{q'}(K)}.
\]
Then, the bound
\[
\|u_\psi\|_{L^{q'}(\Omega_0 \setminus B(0,R))} \le C_K \|\psi\|_{L^{q'}(K)}
\]
is deduced from the explicit representation of $u_\psi$ with the fundamental solution of the Stokes operator in a half-space: as $u_\psi(x) = \int_{K} G(x,y) \psi(y) dy$ and $|G(x,y)| \lesssim C/|x-y|$ (see the appendix), we~find
\[
\|u_\psi\|_{L^{q'}(\Omega_0 \setminus B(0,R))} \lesssim \Bigl\|x \rightarrow \frac{1}{|x|+1}\Bigr\|_{L^{q'}(\Omega_0)}\, \|\psi\|_{L^1(K)} \le C_K \|\psi\|_{L^{q'}(K)}
\]
(here we have used that $q' > 3$). This concludes the proof of \eqref{reg_u_psi}.

Now,
\begin{align*}
2\int_{\Omega^0} D(v): D (u_\psi) &= 2\int_{\Omega^0\setminus \oldbigcup_i {B}_i} D(v): D (u_\psi) - \sum_i \int_{B_i} D(u_{\app}): D(u_\psi)\\
&=\sum_i \int_{\partial B_i} \Sigma(v,q)n \cdot u_\psi d \sigma - \sum_i \int_{B_i} D(u_{\app}): D(u_\psi)\\
&=\sum_i \int_{\partial B_i} \Sigma(v,q)n \cdot \biggl(u_\psi - \dashint u_\psi \biggr)d \sigma - \sum_i \int_{B_i} D(u_{\app}): D(u_\psi)\\
&= \sum_i \int_{A_i} D(v):D \overline{u_{\psi,i}} - \sum_i \int_{B_i} D(u_{\app}): D(u_\psi),
\end{align*}
where, using Lemma \ref{lemme:extension}, $\overline{u_{\psi,i}} \in W^{1,q'}_0(A_i)$ is a divergence-free lifting of $ \overline{u_{\psi,i}}=u_\psi -\nobreak \dashint_{B_i} u_\psi$ on $B_i$, where $A_i =B(X_i,\theta r)\setminus \overline{B(X_i,r)}$ is such that
\begin{equation}\label{eq:bound_u_psi_i}
\|\na \overline{u_{\psi,i}}\|_{L^{q'}(A_i)}\leq C_\theta \|\nabla u_\psi \|_{L^{q'}(B_i)}.
\end{equation}
Note that the sets $A_i$ are disjoint thanks to \eqref{hyp_dmin}. Using the well-known estimate \cite[Prop.\,2.1]{Gerard-VaretMecherbet20}
\[
\|D(v)\|_{L^2(\Omega^0)} \le C\|D(v)\|_{L^2(\oldbigcup_i B_i)} = \|D\left(u_{\app}\right)\|_{L^2(\oldbigcup_i B_i)},
\]
we get, using Hölder inequality, Sobolev embedding $W^{1,q'}(\Omega_0) \subset L^\infty(\Omega_0) $, $q'>3$, together with estimate \eqref{reg_u_psi},
\begin{align*}
\left |\int_K v \cdot \psi \right|& \leq C\,\|D \left(u_{\app}\right)\|_{L^2(\oldbigcup_i B_i)}\,\phi^{1/2}\, \|D \left(u_{\psi}\right)\|_{L^\infty(\oldbigcup_i B_i)}\\
&\leq C_K \| D\left(u_{\app}\right) \|_{L^2(\oldbigcup_i B_i)}\, \phi^{1/2} \|\psi\|_{L^{q'}(K)}\\
&\leq C_K \phi\| D\left(u_{\app} \right)\|_{L^\infty(\oldbigcup_i B_i)}\, \|\psi\|_{L^{q'}(K)}.
\end{align*}
We conclude by observing that on $B_i$ one has
\[
D\left(u_{\app}\right)= \sum_{j \neq i} D\left(\mathcal{U}^{\st}_R[F] (x,X_j)\right)+ D(u_{f,N}).
\]
For the second term at the right-hand side, we~use standard Sobolev embedding and Stokes estimates. For any $p > 3$:
\begin{align*}
\|D\left(u_{f,N}\right)\|_{L^\infty(K)} &\le C \|D\left(u_{f,N}\right)\|_{W^{1,p}(K)}\\
&\le C' \|f (1-1_{\oldbigcup B_i})\|_{L^p(\R^3)} \le C' (\|f\|_{L^\infty} + \|f\|_{L^1}).
\end{align*}
For the first term at the right-hand side we use Proposition \ref{prop:stokeslet} which yields
\begin{align*}
\Bigl\| \sum_{j \neq i} D\left(\mathcal{U}^{\st}_R[F] (\cdot,X_j)\right) \Bigr\|_{L^\infty(B_i)}& \leq \frac{1}{N} \sum_{j \neq i} \frac{C_\theta }{|X_i-X_j|^2} \\
& \leq \frac{C}{N}\sum_{j \neq i}\dashint_{B(X_j,d_{\min}/4)} \frac{1}{|X_i-y|^2}dy \\
&\leq \int_{{}^\mathrm{c}\! B(X_i,d_{\min}/4)} \frac{f^N(y)dy}{|X_i-y|^2}\\
&\leq C(\|f^N\|_\infty+\|f^N\|_1)\leq C,
\end{align*}
with $f^N=\frac{1}{N} \sum_i \sfrac{1_{B(X_i,d_{\min}/4)}}{|B(X_i,d_{\min}/4|}$ such that $\|f^N\|_1=1$ and $\|f^N\|_\infty \leq \sfrac{C}{N d_{\min}^3}$, which is bounded thanks to \eqref{hyp_dmin}.
\begin{rem}\label{discussion_d_min}
More generally, any assumption ensuring
\[
\frac{1}{N} \sum_{j \neq i} \frac{1 }{|X_i-X_j|^2} \leq C
\]
is sufficient. This can be ensured for instance by assuming the following bound on the infinite Wasserstein distance
\[
W_\infty(\rho^N,\rho) \leq CN^{-1/3}
\]
and the constraint
\[
\frac{1}{\sqrt{N}} \lesssim d_{\min}.
\]
Indeed one has
\begin{align*}
\frac{1}{N}\sum_{j \neq i} \frac{1}{|X_i-X_j|^2}&= \int_{Tx\neq X_i}\frac{1}{|Tx-X_i|^2} \,\rho(x) dx,
\end{align*}
with $T$ being the optimal transport plan for the $W_\infty:=W_\infty(\rho^N,\rho)$ Wasserstein distance, $\rho^N=T\# \rho$. We split the integral into two parts:
\begin{align*}
E_1&=\{ x \in \supp \rho,\, Tx \neq X_i,\, |x-X_i| \leq 2W_\infty\},\,
\\
E_2 &= \{ x \in \supp \rho,\, Tx \neq X_i,\, |x-X_i|>2W_\infty\}.
\end{align*}
Hence on $E_1$ one has
\[
\int_{E_1}\frac{1}{|Tx-X_i|^2}\, \rho(x) dx \leq C\,\frac{W_\infty^3}{d_{\min}^2}.
\]
On $E_2$ we have $|Tx-X_i|\geq |x-X_i|-|Tx-x|\geq |x-X_i|/2 $ since $ |Tx-x| \leq W_\infty$, hence
\[
\int_{E_2}\frac{1}{|Tx-X_i|^2}\, \rho(x) dx \leq \int \frac{1}{|x-X_i|^2}\, \rho(x) dx \leq C(\|\rho\|_\infty+\|\rho\|_1),
\]
which shows that
\[
\frac{1}{N}\sum_{j \neq i} \frac{1}{|X_i-X_j|^2} \leq C\Bigl(1+\frac{W_\infty^3}{d_{\min}^2}\Bigr),
\]
see also \cite[Lem.\,2.3]{HoferSchubert23} for more details.\qedhere
\end{rem}
\end{proof}

\subsection{From $u_{\app}$ to $v^N$}
We introduce the intermediate velocity
\begin{equation} \label{eq:Stokes_v^N}
\begin{aligned}
-\Delta v^N + \na q^N & =-\frac{1}{N} \sum_i\delta_{X_i} e+f(1-1_{\oldbigcup_i B_i}) \quad \text{ in } \: \Omega^0, \\
\div v^N & = 0 \quad \text{ in } \: \Omega^0, \\
v^N \vert_{x_1=0} & = 0,
\end{aligned}
\end{equation}
which writes
\[
v^N=-\frac{1}{N}\sum_iG(\cdot,X_i) e+u_{f,N}.
\]
We again use the duality argument to get
\begin{prop}
Let $K \Subset \overline{\Omega^0}$. We have, for any $q<3/2$,
\[
\| u_{\app}- v^N\|_{L^q(K)} \leq C_K R.
\]
\end{prop}
\begin{proof}
Given $K \Subset \overline{\Omega^0}$ and $\psi \in C^\infty_c(K)$, we~again set $u_\psi \in W^{2,q'}(\Omega^0) \cap W^{1,q'}_0(\Omega^0)$ the unique solution to the Stokes equation $-\Delta u_\psi+\nabla q_\psi=\psi $, $\div u_\psi=0$ on $\Omega^0$ with vanishing Dirichlet boundary condition on $x_1=0$ satisfying
\begin{equation}\label{eq:estimation_dualite}
\|u_\psi \|_{W^{2,q'}(\Omega^0)} \leq C_K \|\psi\|_{L^{q'}(K)}.
\end{equation}
We have, using \eqref{eq:Stokes_v^N} and the same notation for $u_\psi$, $\overline{u_{\psi,i}}$ as in the previous proof,\vspace*{-3pt}
\begin{align*}
\int_K(u_{\app}&-v^N)\cdot \psi =\int_{\Omega^0} D(u_{\app}-v^N): Du_{\psi}\\
&=\sum_i \int_{\Omega^0\setminus B_i } D(\mathcal{U}^{\st}_R[F](\cdot,X_i)) : D(u_{\psi}) + \sum_i\frac{1}{N}\,e \cdot u_\psi(X_i)\\
&= \sum_i \int_{\partial B_i}\sigma \left(\mathcal{U}^{\st}_R[F](\cdot,X_i),\mathcal{P}^{\st}_{r}[F](\cdot,y_0)\right) n \cdot u_{\psi} +\sum_i \frac{1}{N}\,e \cdot u_\psi(X_i)\\
&= \sum_i \int_{\partial B_i}\sigma \left(\mathcal{U}^{\st}_R[F](\cdot,X_i), \mathcal{P}^{\st}_{r}[F](\cdot,y_0)\right) n \cdot \left(u_{\psi}- \dashint_{B_i} u_\psi\right) \\
&\hspace*{4.2cm}+ \sum_i\frac{1}{N}\, e \cdot \left(u_\psi(X_i)-\dashint_{B_i} u_\psi \right)\\
&=\sum_i \int_{A_i} D \mathcal{U}^{\st}_R[F](\cdot,X_i): D\overline{u_{\psi,i}} + \sum_i\frac{1}{N}\, e \cdot \left(u_\psi(X_i)-\dashint_{B_i} u_\psi \right)\\
& \leq C \sum_i \left \| D \mathcal{U}^{\st}_R[F](\cdot,X_i) \right\|_{L^2(A_i)} |B_i|^{1/2} \|\nabla u_\psi\|_{L^\infty(K)}+ CR \|\nabla u_\psi\|_{L^\infty(K)}\\
& \leq CR \|\psi\|_{L^{q'}(K)},
\end{align*}
where we used for the first term in the right hand side Hölder inequality, estimate~\eqref{eq:bound_u_psi_i} together with the bound\vspace*{-3pt}
\[
\|\nabla u_\psi\|_{L^\infty(K)} \le C \|\nabla u_\psi\|_{W^{1,q'}(K)} \le C' \|\psi\|_{L^{q'}(K)}
\]
and the energy bound \eqref{eq:energy} for $|F|=1/N$.
\end{proof}

\subsection{From $v^N$ to $u^\varepsilon$}
Keeping in mind that we can restrict to $\eps = 0$, the last step in the proof of Theorem \ref{thm1} is to establish:
\begin{prop}
Let $K \Subset \overline{\Omega^0}$. We have, for any $q<3/2$,\vspace*{-3pt}
\[
\|v^N-u^0\|_{L^q(K)} \leq C_{K} \big(\|{\rho}^N-\rho\|_{(W^{2,q'}(\Omega^0))^*} + \phi \big),
\]
where $u^0$ is the solution of\vspace*{-3pt}
\begin{equation*}
\begin{aligned}
-\Delta u^0 + \na p^0 & = f - \rho e \quad \text{ in } \: \Omega^0, \\
\div u^0 & = 0 \quad \text{ in } \: \Omega^0,\\
u^0 \vert_{\pa \Omega^\eps} & = 0.
\end{aligned}
\end{equation*}
\end{prop}
\begin{proof}
Using again the duality argument we have, for any $\psi \in C^\infty_c(K)$,\vspace*{-3pt}
\begin{align*}
\int_K (v^N-u^\varepsilon)&{}\cdot \psi = \int_{\Omega^0} \nabla (v^N-u^\varepsilon): \nabla u_\psi \\
&= \int_{\Omega^0} ({\rho}^N-\rho) g \cdot u_\psi - \int_{\oldbigcup_i B_i} f \cdot u_\psi \\
&\leq g \|{\rho}^N-\rho\|_{(W^{2,q'}(\Omega^0))^*} \| u_\psi \|_{W^{2,q'}(\Omega^0)} + \sum_i |B_i| \|f\|_{L^\infty(K)} \|\| u_\psi \|_{L^\infty(K)}\\
&\leq C_{K} \big(\|{\rho}^N-\rho\|_{(W^{2,q'}(\Omega^0))^*} + \phi \big) \|\psi\|_{L^{q'}(K)}.
\end{align*}
\begin{rem} \label{rem:Wasserstein_versus_sob_neg}
Since the Wasserstein distance $W_1$ corresponds to the dual norm on Lipschitz functions, one could also use the following estimate
\[
\int_{\Omega^0} ({\rho}^N-\rho) g \cdot u_\psi \leq C W_1(\rho^N,\rho) \|u_\psi\|_{\Lip},
\]
assuming $\rho$ to be a probability measure. However the norm $\| \cdot \|_{(W^{2,q'}(\Omega^0))^*} $ is more accurate since $ W^{2,q'}(\Omega^0)\subset \Lip(\Omega^0)$.\qedhere
\end{rem}
\end{proof}
\section{The approximate continuous model} \label{sec_wall_law}The purpose of this section is to analyze system \eqref{limit_equation}, and to prove Theorem \ref{thm2} and Corollary \ref{coro1}. We assume here that $f$ and $\rho$ are smooth and compactly supported in~$\overline{\Omega^0}$.

\subsection{Asymptotic expansion}
We will construct an approximation of the solution $(u^\eps,p^\eps)$ of \eqref{limit_equation}, of the following form:
\begin{equation*}
\left\{
\begin{aligned}
(u^\eps_{\app},p^\eps_{\app})(x_1,x_2,x_3) & \approx \sum_{i=0}^m \eps^i (u^i, p^i)(x_1,x_2,x_3), \quad x = (x_1,x_2,x_3) \in \Omega^0, \\
(u^\eps_{\app}, p^\eps_{\app}(x_1,x_2,x_3) & \approx \sum_{i=0}^m \eps^i (U^i, P^i)(x_1/\eps,x_2,x_3), \quad x = (x_1,x_2,x_3) \in \Omega^\eps\setminus\Omega^0.
\end{aligned}
\right.
\end{equation*}
Note that this approximation is made of two parts: an interior one, in $\Omega^0$, with a regular expansion in $\eps$, and a boundary layer part, localized in the depletion layer $\Omega^\eps \setminus\nobreak \Omega^0$. The boundary layer expansion involves boundary layer profiles $U^i = U^i(s,x_2,x_3)$, with the variable $s \in (-1,0)$ that stands for
$x_1/\eps$. It is also convenient to set:
\[
(u^i, p^i)= 0, \quad (U^i, P^i) = 0 \quad \text{ for } \: i < 0.
\]
Plugging the interior expansion in the Stokes equation, we~find: for all $i \ge 0$, in $\Omega^0$
\begin{equation} \label{stokesint}
\begin{aligned}
-\Delta u^i + \na p^i & = \delta_{0i} \big(f - \rho e \big), \\
\div u^i & = 0.
\end{aligned}
\end{equation}
Plugging the boundary layer expansion in the Stokes equation, we~find: for all $i \ge 0$, for all $(s,x_2,x_3) \in (-1,0) \times \R^2$:
\begin{align}
-\pa^2_s U^i_1 - (\pa^2_2 + \pa^2_3) U^{i-2}_1 + \pa_s P^{i-1} & = 0, \label{stokesx} \\
-\pa^2_s U^i_2 - (\pa^2_2 + \pa^2_3) U^{i-2}_2 + \pa_2 P^{i-2} & = 0, \label{stokesy}\\
-\pa^2_s U^i_3 - (\pa^2_2 + \pa^2_3) U^{i-2}_3 + \pa_3 P^{i-2} & = 0, \label{stokesz} \\
\pa_s U^i_1 + \pa_2 U^{i-1}_2 + \pa_3 U^{i-1}_3 & = 0 \label{stokesdiv}
\end{align}
The Dirichlet condition $u^\eps_{\app}\vert_{\pa \Omega^\eps} \approx 0$ further yields: for all $i \ge 0$,
\begin{equation} \label{dirichlet}
U^i\vert_{s = -1} = 0.
\end{equation}
Eventually, conditions
\[
[u^\eps_{\app}]\vert_{\pa \Omega^0} \approx 0, \quad [\Sigma(u^\eps_{\app},p^\eps_{\app})n]\vert_{\pa \Omega^0} \approx 0,
\]
which reflect continuity of the velocity field and the stress tensor at the interface $\pa \Omega^0$ give for all $i \ge 0$:
\begin{align}
U^i\vert_{s=0} & = u^i\vert_{x_1=0} \label{continuity_vel} \\
\pa_s U^i \vert_{s=0} - P^{i-1}\vert_{s=0} (1,0,0)^t & = \pa_1 u^{i-1}\vert_{x_1=0} - p^{i-1}\vert_{x_1=0} (1,0,0)^t. \label{continuity_stress}
\end{align}

\subsubsection*{Computation of the first terms.} We compute $U^0, u^0,p^0$. First, from \eqref{stokesdiv} and \eqref{dirichlet}, we~get $\pa_s U^0_1 = 0$ and $U^0_1\vert_{s=-1} = 0$, which implies $U^0_1 = 0$. Then, from \eqref{stokesy}, \eqref{dirichlet} and~\eqref{continuity_stress}, we~find
\[
\pa^2_s U^0_2 = 0, \quad U^0_2\vert_{s=-1} = 0, \quad \pa_s U^0_2\vert_{s=0} = 0,
\]
which leads to $U^0_2 = 0$. Similarly, $U^0_3 = 0$. Hence, $U^0 = 0$. Next, we~consider \eqref{stokesint} and~\eqref{continuity_vel} at rank $i=0$. They imply that $(u^0,p^0)$ satisfies \eqref{limit_equation}. In~particular, as~$f$ belongs to $L^{6/5}(\Omega_0) \cap H^m(\Omega^0)$, $(\na u^0,p^0)$ belongs to $H^m(\Omega^0)$ for all $m$.

\subsubsection*{Computation of next order terms.} We now assume that $i \ge 1$ and that profiles $U^k$, $P^{k-1}$, $u^k$, $p^k$ are known for all $k \le i-1$, with:
\[
U^k, P^{k-1} \in H^\infty((-1,0) \times \R^2), \quad \na u^k, p^k \in H^\infty(\Omega^0).
\]
We now show how to construct $U^i,P^{i-1}, u^i, p^i$. Considering \eqref{stokesdiv} and \eqref{dirichlet} yields
\[
U^i_1(s,\cdot) = - \int_{-1}^s (\pa_2 U^{i-1}_2 + \pa_3 U^{i-1}_3)(t,\cdot) dt \in H^\infty((-1,0) \times \R^2).
\]
Next, we~consider \eqref{stokesx} and \eqref{continuity_stress} which allow the calculation of \hbox{$P^{i-1} \!\!\in\! H^\infty(\!({-}1,0)\!\times\! \R^2)$}:
\[
P^{i-1}(s,\cdot) = \pa_s U^i_1\vert_{s=0} - \pa_1 u^{i-1}_1\vert_{x_1=0} + p^{i-1}\vert_{x_1=0}
+ \int_0^s \big(\pa^2_s U^i_1 + (\pa^2_2 + \pa^2_3) U^{i-2}_1 \big)(t,\cdot) dt.
\]
Next, we~consider \eqref{stokesy}--\eqref{dirichlet}--\eqref{continuity_stress}: these relations imply
\[
U^i_2(s,\cdot) = - \int_{-1}^s \int_{0}^t \Big((\pa^2_2 + \pa^2_3) U^{i-2}_2 + \pa_2 P^{i-2} \Big)(t',\cdot) dt' dt + (s+1) \pa_1 u^{i-1}_2\vert_{x_1=0}.
\]
A similar expression holds for $U^i_3$. Both $U^i_2, U^i_3 \in H^\infty((-1,0) \times \R^2)$. Eventually, we~see by \eqref{stokesint} and \eqref{continuity_vel} that $(u^i,p^i)$ solve a homogeneous Stokes equation with the inhomogeneous Dirichlet condition
\[
u^i\vert_{x_1=0} = U^i\vert_{x_1=0}.
\]
By standard regularity results for the Stokes equation, \cf
\cite[Th.\,IV.3.2 \& IV.3.3]{Galdi}, $\na u^i, p^i \in H^\infty(\Omega^0)$.

\subsubsection*{Equations for $u^1$ and $u^2$.}
Anticipating the discussion in Section \ref{sec_final_comments}, it is worth specifying the systems satisfied by $u^1,p^1$ and $u^2,p^2$. From above equations, we~deduce
\[
U^1_1 = 0, \quad P^0 = -\pa_1 u^0_1\vert_{x_1=0} + p^0\vert_{x_1=0}, \quad (U^1_2, U^1_3) = (s+1) (\pa_1 u^0_2,\pa_1 u^0_3)\vert_{x_1=0}.
\]
Hence,
\begin{equation} \label{systemu1}
\begin{aligned}
-\Delta u^1 + \na p^1 & = 0, \\
\div u^1 & = 0, \\
u^1\vert_{x_1=0} & = (0,\pa_1 u^0_2, \pa_1 u^0_3)\vert_{x_1=0}.
\end{aligned}
\end{equation}
Then,
\begin{align*}
U^2_1(s,\cdot) &= - \int_{-1}^s (\pa_2 U^{1}_2 + \pa_3 U^{1}_3)(t,\cdot) dt \\
& = - \int_{-1}^s (t+1)(\pa_2 \pa_1 u^0_2 + \pa_3 \pa_1 u^0_3)\vert_{x_1=0}dt = \frac{(s+1)^2}{2} \pa^2_1 u^0_1\vert_{x_1=0},
\end{align*}
while
\begin{align*}
(U^2_2, U^2_3)&= (s+1) (\pa_1 u^1_2, \pa_1 u^1_3)\vert_{x_1=0} + \frac{1}{2} (s-1) (s+1) (\pa_2 P^0, \pa_3 P^0) \\
& = (s+1) (\pa_1 u^1_2, \pa_1 u^1_3)\vert_{x_1=0}\\
&\hspace*{2cm} + \frac{1}{2} (s-1) (s+1) \big(-\pa_1\pa_2 u^0_1 + \pa_2 p^0, -\pa_1\pa_3 u^0_1 + \pa_3 p^0)\vert_{x_1=0}.
\end{align*}
We recover the system
\begin{equation} \label{systemu2}
\begin{aligned}
-\Delta u^2 + \na p^2 & = 0, \\
\div u^2 & = 0, \\
u^2_1\vert_{x_1=0} & = \frac{1}{2}\pa^2_1 u^0_1\vert_{x_1=0}, \\
u^2_2\vert_{x_1=0} & = \pa_1 u^1_2\vert_{x_1=0} + \frac{1}{2}(\pa_1\pa_2 u^0_1 - \pa_2 p^0)\vert_{x_1=0},\\
u^2_3\vert_{x_1=0} & = \pa_1 u^1_3\vert_{x_1=0} + \frac{1}{2}(\pa_1\pa_3 u^0_1 - \pa_3 p^0)\vert_{x_1=0}.
\end{aligned}
\end{equation}

\subsection{Proof of Theorem \ref{thm2}}
We will show here that the expansion constructed in the previous section provides an approximate solution of the true solution $u^\eps$ of \eqref{limit_equation}. Given $m \in \N$, we~first introduce $\tilde{u}^{m+2}$ the solution in $\dot{H}^1(\Omega^\eps)$ of the Stokes problem
\begin{equation*}
\left\{
\begin{aligned}
-\Delta \tilde{u}^{m+2} + \na \tilde{p}^{m+2} & = 0 \quad \text{in } \Omega^\eps, \\
\div \tilde{u}^{m+2} & = - \pa_1 u^{m+2}_1 1_{\Omega^0} \quad \text{in } \Omega^\eps, \\
\tilde{u}^{m+2}\vert_{\pa \Omega^\eps} & = 0.
\end{aligned}
\right.
\end{equation*}
Such $\tilde{u}^{m+2}$ exists, \cf \cite[Th.\,IV.3.3]{Galdi}, and satisfies
\[
\|\na \tilde{u}^{m+2}\|_{L^2(\Omega^\eps)} \lesssim \|\pa_1 u^{m+2}_1 1_{\Omega^0} \|_{L^2(\Omega^\eps)} \lesssim 1.
\]
We then define $(u^\eps_{\app}, p^{\eps}_{\app})$ as follows:
\begin{equation*}
\begin{aligned}
(u^\eps_{\app},p^\eps_{\app})(x,y,z) & = \sum_{i=0}^{m+1} \eps^i (u^i, p^i)(x,y,z) + \eps^{m+2} (u^{m+2}_1,0,0,0)(x,y,z) \\
& \hspace*{1cm}+ \eps^{m+2} (\tilde{u}^{m+2},\tilde{p}^{m+2})(x,y,z), \qquad X = (x,y,z) \in \Omega^0, \\
(u^\eps_{\app}, p^\eps_{\app})(x,y,z) & = \sum_{i=0}^{m+1} \eps^i (U^i, P^i)(x/\eps,y,z) + \eps^{m+2} (U^{m+2}_1,0,0,0)(x/\eps,y,z) \\
& \hspace*{1cm}+ \eps^{m+2} (\tilde{u}^{m+2},\tilde{p}^{m+2})(x,y,z), \qquad \: X = (x,y,z) \in \Omega^\eps\setminus\Omega^0.
\end{aligned}
\end{equation*}
We then introduce $v^\eps = u^\eps - u^\eps_{\app}$, $q^\eps = p^\eps - p^\eps_{\app}$, solving
\begin{align}
-\Delta v^\eps + \na q^\eps & = - \eps^{m+2} (\Delta u^{m+2}_1,0,0) \quad \text{ in } \: \Omega^0,\notag \\
-\Delta v^\eps + \na q^\eps & = - \sum_{i=m}^{m+1} \eps^i (\pa^2_2 + \pa^2_3) (\eps U^{i+1}_1,U^i_2,U^i_3) - \sum_{i=m}^{m+1} \eps^i (0, \pa_2 P^i, \pa_3 P^i) \quad \text{ in } \: D^\eps, \notag \\
\div v^\eps & = 0 \quad \text{ in } \: \Omega^\eps, \\
[v^\eps]\vert_{\pa \Omega^0} & = 0, \notag\\
[\Sigma(v^\eps,q^\eps)n]\vert_{\pa \Omega^0} & = - \eps^{m+1}(\eps \pa_1 u^{m+2}_1\vert_{x_1=0},\pa_1 u^{m+1}_2\vert_{x_1=0}, \pa_1 u^{m+1}_3\vert_{x_1=0}), \notag \\
v^\eps \vert_{\pa \Omega^\eps} & = 0. \notag
\end{align}
Performing a simple energy estimate, we~find
\begin{align*}
\int_{\Omega^\eps} |\na v^\eps|^2 = \eps^{m+2} \int_{\Omega^0} \na u_1^{m+2} : \na v_1^\eps - \int_{D^\eps} R^\eps \cdot v^\eps + \int_{\pa \Omega^0} r^\eps \cdot v^\eps,
\end{align*}
where
\begin{align*}
R^\eps & = - \sum_{i=m}^{m+1} \eps^i (\pa^2_2 + \pa^2_3) (\eps U^{i+1}_1,U^i_2,U^i_3) - \sum_{i=m}^{m+1} \eps^i (0, \pa_2 P^i, \pa_3 P^i) = O(\eps^m), \\
r^\eps & = \eps^{m+1}(\eps \pa_1 u^{m+2}_1\vert_{x_1=0},\pa_1 u^{m+1}_2\vert_{x_1=0}, \pa_1 u^{m+1}_3\vert_{x_1=0}) = O(\eps^{m+1}).
\end{align*}
By the Poincaré and trace inequalities
\[
\| v^\eps \|_{L^2(D^\eps)} \lesssim \eps \|\na v^\eps \|_{L^2(D^\eps)}, \quad \| v^\eps \|_{L^2(\pa\Omega^0)} \lesssim \sqrt{\eps} \|\na v^\eps \|_{L^2(D^\eps)},
\]
we infer easily
\[
\|\na v^\eps \|_{L^2(\Omega^\eps)} \lesssim \eps^{m+1}.
\]
Restricting to $\Omega^0$, Theorem \ref{thm2} follows.

\subsection{Wall law of Navier type : proof of Corollary \ref{coro1}}
We conclude Section \ref{sec_wall_law} by proving Corollary \ref{coro1}.
We know from \eqref{thm2} that
\[
u^\eps = u^0 + \eps u^1 + O(\eps^2) \quad \text{ in } \: \dot{H}^1(\Omega^0).
\]
Hence, Corollary \ref{coro1} will follow from the estimate
\[
\|\na (u^S - u^S_{\app})\|_{L^2(\Omega^0)} = O(\eps^2), \quad u^S_{\app} = u^0 + \eps u^1.
\]
The difference $v^{S} = u^S - u^S_{\app}$ satisfies
\begin{equation*}
\begin{aligned}
-\Delta v^S + \na q^S & = 0 \quad \text{ in } \: \Omega^0, \\
\div v^S & = 0 \quad \text{ in } \: \Omega^0, \\
v^S\vert_{\pa \Omega^0} & = \eps (0, \pa_1 v^S_2, \pa_1 v^S_3)\vert_{\pa \Omega^0} + \eps^2 (0, \pa_1 u^1_2, \pa_1 u^1_3)\vert_{\pa \Omega^0}.
\end{aligned}
\end{equation*}
Let $V$ be the solution of the Stokes problem:
\begin{equation*}
\begin{aligned}
-\Delta V + \na Q & = 0 \quad \text{ in } \: \Omega^0, \\
\div V & = 0 \quad \text{ in } \: \Omega^0, \\
V\vert_{\pa \Omega^0} & = (0, \pa_1 u^1_2, \pa_1 u^1_3)\vert_{\pa \Omega^0}.
\end{aligned}
\end{equation*}
Then, $\na V$ belongs to $H^\infty(\Omega^0)$. Moreover, the function $w^S = v^S - \eps^2 V$ satisfies
\begin{equation*}
\begin{aligned}
-\Delta w^S + \na r^S & = 0 \quad \text{ in } \: \Omega^0, \\
\div w^S & = 0 \quad \text{ in } \: \Omega^0, \\
w^S\vert_{\pa \Omega^0} & = \eps (0, \pa_1 w^S_2, \pa_1 w^S_3)\vert_{\pa \Omega^0} + \eps^3 (0, \pa_1 V_2, \pa_1 V_3)\vert_{\pa \Omega^0}.
\end{aligned}
\end{equation*}
The standard estimate
\begin{equation*}
\|\na w^S\|_{L^2(\Omega^0)}^2 + \frac{1}{\eps} \int_{\pa \Omega^0}
(|w^S_2|^2 + |w^S_3|^2) = \eps^2 \int_{\pa \Omega^0} \big(\pa_1 V_2 \, w^S_2 \: + \: \pa_1 V_3 \, w^S_3\big)
\end{equation*}
combined with Young's inequality yields:
\begin{equation*}
\|\na w^S\|_{L^2(\Omega^0)}^2 + \frac{1}{2\eps} \int_{\pa \Omega^0} (|w^S_2|^2 + |w^S_3|^2) \le \frac{1}{2}\eps^5 \int_{\pa \Omega^0} \big(|\pa_1 V_2|^2 + \, |\pa_1 V_3|^2 \big).
\end{equation*}
It follows that
$\|\na w^S\|_{L^2(\Omega^0)} = O(\eps^{5/2})$ and finally $\|\na v^S\|_{L^2(\Omega^0)} = O(\eps^2)$. The result follows.

\section{Final comments} \label{sec_final_comments}

\subsection{Apparent slip} \label{subsec_apparent}
By Corollary \ref{coro1}, the approximation of system \eqref{limit_equation} by system~\eqref{improved_approximation} is more accurate than the one by \eqref{crude_approx}, where the effect of the depletion layer is not taken into account and homogeneous boundary conditions are considered. The boundary condition in \eqref{improved_approximation} is called a (Navier) slip boundary condition. To understand this terminology, one can consider the case of a shear flow driven by a constant downward pressure gradient:
\[
- \Delta u + \na p = - e, \quad \div u = 0 \quad \text{ in } (0,1) \times \R.
\]
In the case of Dirichlet conditions $u\vert_{x_1=0} = u\vert_{x_1=1} = 0$, the solution is the usual Poiseuille flow
\[
u^0(x_1,x_2) = \big(0, \tfrac{1}{2} x_1 (x_1-1)\big),
\]
while in the case of Navier conditions
\[
u\vert_{x_1=0} = (0, \pa_1 u_2\vert_{x_1=0}), \quad u\vert_{x_1=1} = (0, -\pa_1 u_2\vert_{x_1=1}),
\]
we find
\[
u^S(x_1,x_2) = \big(0,\tfrac{1}{2} x_1 (x_1-1) - \tfrac12\big),
\]
which has a non-zero downward component at the boundary. This corresponds to a phenomenon of {\em apparent slip}, at the artificial boundary $\pa \Omega^0$ (while real no-slip takes place at the rigid wall $\pa \Omega^\eps$). As mentioned in this introduction, this phenomenon of apparent slip was noticed in various experiments on sheared suspensions. It is here validated mathematically by Corollary \ref{coro1}.

Let us stress that this improved approximation is established in Section \ref{sec_wall_law} starting from the continuous model \eqref{limit_equation}. Actually, as we are interested in the ``real'' suspension, governed by \eqref{main_equation}, the point would rather be to show that for some norm
\begin{equation} \label{improved_real}
\|u^{N,\eps} - u^S\| < \|u^{N,\eps} - u^0\|.
\end{equation}
By Theorem \ref{thm1}, for all $q < 3/2$, $K \Subset \overline{\Omega_0}$,
\begin{align*} \|u^{N,\eps} - u^0\|_{L^q(K)} &\ge \|u^\eps - u^0\|_{L^q(K)} - \|u^{N,\eps} - u^\eps\|_{L^q(K)}\\
&\ge \|u^\eps - u^0\|_{L^q(K)} - C (R+\phi + \|\rho^N-\rho\|_{(W^{2,q'}(\Omega^0))^*}).
\end{align*}
If the solution $u^1$ of \eqref{systemu1} is non-identically zero on $K$, then by Theorem \ref{thm2}: \hbox{$\|u^\eps - u^0\|_{L^q(K)} \approx \eps$}, so that for a large enough constant $C_0$ (independent of $N$, $R$ and of course $\eps$),
\[
\|u^{N,\eps} - u^0\|_{L^q(K)} \ge C_0 \, \eps - C (R+\phi + \|\rho^N-\rho\|_{(W^{2,q'}(\Omega^0))^*}).
\]
On the other hand, combining Theorems \ref{thm1} and \ref{thm2}, we~also have for some constant~$C_1$
\begin{align*}
\|u^{N,\eps} - u^S\|_{L^q(K)} &\le \|u^\eps - u^S\|_{L^q(K)} + \|u^{N,\eps} - u^\eps\|_{L^q(K)}\\
&\le C_1\,\eps^2 + C (R+\phi + \|\rho^N-\rho\|_{(W^{2,q'}(\Omega^0))^*}).
\end{align*}
This ensures that condition \eqref{improved_real} is satisfied as soon as
\begin{equation} \label{stronger_assumption_size}
\frac{C_0}{2 C_1} > \eps \ge \mathcal{C} \big(R+ \phi + \|\rho^N-\rho\|_{(W^{2,q'}(\Omega^0))^*} \big), \quad \mathcal{C} := \frac{4C}{C_0}.
\end{equation}
The lower bound on $\eps$ in this inequality is of course more stringent than the simple condition \eqref{H1}. However, we~believe it is still relevant for the following reasons:
\begin{itemize}
\item There is experimental evidence that some suspensions have a depletion layer obeying
\begin{equation} \label{5'} \tag{5'}
\eps \geq \mathcal{C'} R
\end{equation}
for some constant $\mathcal{C}' $ much larger than $1$. For instance, \cite{Koponen} considers suspensions of microfibrillated cellulose in a pipe flow, for which the depletion layer has size $\approx 200 \mu \mathrm{m}$ whereas the typical length of one particle is $8 \mu \mathrm{m}$, resulting in a constant $\mathcal{C'} = 25$. Although our constant $\mathcal{C}$ in \eqref{stronger_assumption_size} is not very explicit, the requirement $\eps \ge \mathcal{C} R$ contained in \eqref{stronger_assumption_size} is reasonable with regards to this kind of suspensions, and our Corollary \ref{coro1} can be seen as a justification of the apparent slip detected experimentally.
\item In the case of periodic distributions one has $\|\rho^N-\rho\| \sim N^{-2/3}$, see \cite{HillairetHoefer23}, whereas for random particles, drawn independently with the same law, one has roughly $\|\rho^N-\nobreak\rho\| \sim N^{-1/2}$ (due to central limit theorem). In~the first, \resp second setting, under the requirement that
\[
N^{-\sfrac23} \lesssim R \lesssim N^{-\sfrac12}, \quad \text{ \resp } R \sim \frac{1}{\sqrt{N}}
\]
assumption \eqref{stronger_assumption_size} reduces to the realistic hypothesis \eqref{5'}.
\item The constraint coming from the $\phi$ term in \eqref{stronger_assumption_size} could be relaxed, by including some effective viscosity term in the continuous Stokes equation \eqref{limit_equation}. Typically, one can replace $\phi$ by $\phi^2$ if the Einstein's effective viscosity is taken into account, that is replacing \eqref{limit_equation} by
\[
-\div\bigl((1+ \tfrac52 \phi \rho) \na u \bigr) + \na p = f - \rho e, \quad \div u = 0 \quad \text{ in } \Omega^\eps, \quad u\vert_{\pa \Omega^\eps = 0}.
\]
See \cite{Hofer&Schubert} for detailed statements. For periodic or random stationary configurations, one can even have an effective viscosity tensor accurate at any power of $\phi$, relaxing completely the constraint. See \cite{DuerinckxGloria19}. The subsequent derivation of wall laws is not affected by the introduction of such refined Stokes equation.
\end{itemize}

To establish the validity of a slip law of Navier type under the mere assumption~\eqref{H1} would in our opinion require different ideas. In~particular, to use the intermediate system \eqref{limit_equation} would not be valid anymore, because if $\eps$ is comparable to $R$, the microscopic structure of the distribution of particles (for instance the shape of the particles) should be involved. This can be seen in similar problems, like the derivation of wall laws for porous media \cite{Mikelic1996} or for rough boundaries \cite{Mikelic2009,Basson2008}. In~such problems, the approach used is based on homogenization tools, and the slip length appears as an averaged quantity related to a kind of cell problem. To go beyond the periodic setting raises strong mathematical issues. Here, we~have rather chosen to consider pretty general distributions of particles.

\subsection{Intrinsic convection} \label{subsecintrinsic}
In all our analysis, it is implicit that the first terms~$u^0$ and~$u^1$ in our expansion of $u^\eps$ are non-zero. This is very natural in view of our assumptions, where the limit density of the particles distribution $\rho$ is {\em inhomogeneous} (and even compactly supported). This inhomogeneity drives a non-zero global flow with field $u^0$. Such setting excludes the case of a homogeneous suspension, meaning with constant density over
$\Omega^0$ (and no driving force). Nevertheless, we~can still in this degenerate case use the continuous model \eqref{limit_equation}, with $\rho=1$ and $f=0$, and perform at least formally the asymptotic expansion seen in the proof of Theorem \ref{thm2} (an analogue of this theorem for constant $\rho$ would require to change the functional spaces, as there is no decay anymore at infinity).

Such an approach for homogeneous suspensions can be found (although sketchy) in~\cite{BFAH96}, for a 2d channel of the form $(-\eps,1) \times \R$. See also \cite{BFBA98}. This can be straightforwardly adapted to our domain $\Omega^\eps$. First, when $\rho=1$ and $f=0$, the solution of \eqref{crude_approx} is given by
\[
u^0(x) = 0, \quad p^0(x) = -x_3.
\]
It follows from \eqref{systemu1} that $u^1=0, \: p^1 = 0$.
Going one step further and considering \eqref{systemu2}, we~find that $u^2$ satisfies the boundary condition
\[
u^2\vert_{x_1=0}= (0,0,\tfrac{1}{2}).
\]
This gives
\[
u^\eps\vert_{x_1=0} \approx (0,0, \tfrac{1}{2} \eps^2).
\]
Back to dimensional variables, we~find
\[
u^\eps_3\vert_{x_1=0} \approx \frac{\eps^2 mgN}{2 \mu L^3} = \frac{g}{2\mu} \eps^2 (\rho_s-\rho_f) \phi,
\]
using that the mass of a spherical particle (rectified by buoyancy) is \hbox{$m \!=\! \frac{4}{3} \pi R^3 (\rho_p \!-\! \rho_f)$} with $\rho_p$ and $\rho_f$ the particle and fluid masses per unit, while the solid volume fraction is $\phi = \frac{4}{3}\pi R^3/L^3 N$. Introducing the settling speed of a single spherical particle
\[
V_0 = 2R^2 (\rho_p - \rho_f)g/9\mu,
\]
this can be rewritten as
\[
u^\eps_3\vert_{x_1=0} \approx \frac{\eps^2}{R^2} \frac{9}{4} V_0 \phi,
\]
which agrees with the formula found in \cite{BFAH96} for the special case $\eps = R$ considered there.
We see that contrary to the case examined in paragraph \ref{subsec_apparent}, the velocity just outside the depletion layer (at $x_1=0$) is pointing upward. For zero mean flux in a channel, this upward velocity near the layers is compensated by an extra downward velocity in the middle of the channel, as discussed in \cite{BFAH96}. This extra flow is tagged as intrinsic convection in the literature.

Still, we~stress that this formal analysis is conducted at the level of the continuous model \eqref{limit_equation}. If we want to establish rigorously the phenomenon of intrinsic convection, we~need to relate it to the microscopic model \eqref{main_equation}. Arguing as in Section \ref{subsec_apparent}, we~see that it requires
\[
\eps^2 \ge \mathcal{C} \big(R+ \phi + \|\rho^N-\rho\|_{(W^{2,q'}(\Omega^0))^*} \big).
\]
This corresponds to much wider depletion layers than those considered in \cite{BFAH96}. Beyond the question of mathematical justification, this may explain the lack of experimental evidence for this phenomenon.

\appendix
\section*{Appendix. Estimates for the fundamental solution of the Stokes equation on the half space}
\refstepcounter{section}\label{appendix}
Let us first introduce some notations. We denote by $(\Phi,Q)$ the fundamental solution to the Stokes equation on $\R^3$ given by
\[
\Phi(x)=\frac{1}{8\pi}\Bigl(\frac{\mathbb{I}}{|x|}+\frac{x\otimes x}{|x|^3} \Bigr),\quad Q(x)=\frac{1}{4\pi} \frac{x}{|x|^3},
\]
where $\mathbb{I}$ the identity matrix in $3d$. For any $F=(F_1,F_2,F_3) \in \R^3$, we~denote by $F^I=(-F_1,F_2,F_3)$.
We rely on \cite[\S 3.1]{Gimbutas_Greengard_Veerapaneni_2015} for the following result
\begin{prop}\label{prop_appendix}
Let $y \!\in\! (0,+\infty)\times \R^2$ and denote by $G(x,y)$ the unique solution~to
\[
-\Delta_x G+\nabla_x p = \delta_{y} \mathbb{I}, \quad \div_x G= 0, \text{ on } \: (0,+\infty)\times \R^2, \quad G(x,y)\vert_{x_1=0}=0.
\]
\begin{enumerate}
\item We have
\[
G(x,y)=\Phi(x-y)-\Phi^I(x-y^I)-A(x,y),
\]
with
\begin{align*}
A_{i,j} &= (x_1 \partial_{x_i}-\delta_{i1})\Bigl(\frac{1}{4\pi}\frac{(e_j^I)_1}{|x-y^I|}+\frac{y_1}{4\pi} \frac{(x-y^I)\cdot e_j^I}{|x-y^I|^3} \Bigr), \quad 1 \leq i,j\leq 3,\\
(\Phi^I)_{i,j}&=\Phi_{i,j}(1-2\delta_{i1}), \quad 1 \leq i,j\leq 3,
\end{align*}
where $\{e_1,e_1,e_3\}$ stands for the canonical basis of $\R^3$.
\item There exists a constant $C>0$ such that for all $x,y \in (0,+\infty)\times \R^2 $, $x\neq y$
\[
|G(x,y)| \leq \frac{C}{|x-y|}, \quad | \nabla_1 G(x,y)| + | \nabla_y G(x,y)| \leq \frac{C}{|x-y|^2}.
\]
\end{enumerate}
\end{prop}

\skpt
\begin{proof}
\begin{enumerate}
\item The explicit formula of $G$ can be found in the nice paper \cite[Formula (21)]{Gimbutas_Greengard_Veerapaneni_2015}. Let us recall for completeness the main elements of its derivation. Let $F\in \R^3$, the idea is that subtracting to $\Phi(x-y)F$ its reflected point force $\Phi(x-y^I)F^I$ yields a velocity $u$
\[
u(x)=\Phi(x-y)F-\Phi(x-y^I)F^I,
\]
which satisfies $-\Delta u +\nabla p=F \delta_y$, $\div u=0$ on $(0,+\infty)\times\R^2$ with $u_3(x)_{|_{x_1=0}}=u_2(x)_{|_{x_1=0}}=0$ and
\[
u_1(x)_{|_{x_1=0}}=- \frac{1}{4\pi}\frac{F_1^I}{|x-y^I|}-\frac{y_1}{4\pi} \frac{F^I \cdot (x-y^I)}{|x-y^I|^3}.
\]
On the right-hand side, one recognizes the harmonic potential due to a simple charge and the one due to a dipole with orientation $F^I$.
Such an inhomogeneous Dirichlet condition can be corrected by introducing the so called Papkovich-Neuber correction \cite[Def.\,2.1]{Gimbutas_Greengard_Veerapaneni_2015}
\[
u^C(x)=x_1 \nabla \phi-e_1 \phi(x), \quad p^C(x)=2\partial_{x_1}\phi(x),
\]
which satisfies $-\Delta u^C+\nabla p^C =0$, $\div u^C=0$ for any harmonic function $\phi$ and $u^C_{|_{x_1=0}}=-\phi(x)_{|_{x_1=0}}e_1$. Hence in order to correct the boundary condition of $u$ one needs to remove the term $u^C$, with
\[
\phi(\mathbf{x})=\frac{1}{4\pi}\frac{F_1^I}{|x-y^I|}+\frac{y_1}{4\pi} \frac{F^I \cdot (x-y^I)}{|x-y^I|^3}.
\]
This yields the claimed formula.
\item For the decay, we~use the following properties for any $x, y \in (0,+\infty)\times \R^2 $, $x\neq y$,
\[
|x-y^I|\geq |x-y|, \quad \frac{|y_1|}{|x-y^I|} \leq 1, \quad \frac{|x_1|}{|x-y^I|} \leq 1,
\]
together with the following estimates
\begin{align*}
|A(x,y)| &\leq C \Bigl(\frac{|x_1|}{|x-y^I|}+ 1 \Bigr) \Bigl(\frac{1}{|x-y^I|}+ \frac{|y_1|}{|x-y^I|^2} \Bigr),
\\
|\nabla_x A(x,y)|+|\nabla_y A(x,y)| &\leq C \Bigl(\frac{|x_1|}{|x-y^I|}+ 1 \Bigr) \Bigl(\frac{1}{|x-y^I|^2}+ \frac{|y_1|}{|x-y^I|^3} \Bigr).\qedhere
\end{align*}
\end{enumerate}
\end{proof}

\backmatter
\bibliographystyle{jepplain+eid}
\bibliography{gerard-varet-mecherbet}
\end{document}
